\documentclass[11pt]{article}
\usepackage[margin=1in]{geometry}
\usepackage{fontspec}
\usepackage{hyperref}
\usepackage{orcidlink}
\usepackage{fvextra}
\fvset{breaklines=true,breakanywhere=true,fontsize=\small}
\title{Dossier de Registro: padic-grn-viz\\\large Visualización y reportes para análisis p-ádico reproducible}
\author{
Jesús Rogelio Pérez-Buendía\,\orcidlink{0000-0002-7739-4779}
\and
Víctor Nopal-Coello\,\orcidlink{0000-0003-2608-3636}
\and
Yuriria Cortés-Poza\,\orcidlink{0000-0003-1856-9206}
}
\date{2026-04-25}

\begin{document}
\maketitle
\begin{abstract}
\texttt{padic-grn-viz} proporciona herramientas de visualización para comunicar resultados de análisis p-ádico en sistemas dinámicos discretos. Este dossier documenta su propósito, forma de uso, implementación, pruebas y evidencia de ejecución, con enfoque en claridad para personas usuarias y reproducibilidad técnica para registro de software.
\end{abstract}
\tableofcontents
\newpage

\section{Introducción}
Este software transforma salidas numéricas y estructurales de padic-grn-core y padic-grn-search en gráficos y artefactos visuales reutilizables. Su papel dentro de la suite es facilitar interpretación, comparación y comunicación de resultados tanto en contextos de investigación como en reportes institucionales.

La arquitectura actual se centra en motores nativos de visualización y comandos de línea para generar figuras reproducibles a partir de archivos de entrada bien definidos. Con ello, padic-grn-viz mejora la accesibilidad del análisis para personas lectoras no especializadas sin perder rigor técnico.

\section{Datos generales}
\begin{itemize}
  \item Nombre: \texttt{padic-grn-viz}
  \item Versión: \texttt{0.1.0}
  \item Lenguaje: Python 3.10+
  \item Dependencia principal: \texttt{matplotlib}
  \item Tipo: generación de visualizaciones reproducibles
  \item DOI: \texttt{DOI\_SOFTWARE\_VIZ}
\end{itemize}

\subsection{Correos electrónicos oficiales de autores}
\begin{itemize}
  \item Jesús Rogelio Pérez-Buendía: \texttt{rogelio.perez@cimat.mx}
  \item Víctor Nopal-Coello: \texttt{victor\_nopal@uaeh.edu.mx}
  \item Yuriria Cortés-Poza: \texttt{yuriria.cortes@iimas.unam.mx}
\end{itemize}

\subsection{Financiamiento}
\begin{itemize}
  \item Modelos matemáticos y computacionales no convencionales para el estudio y análisis de problemas relevantes en Biología, CF 2019/217367.
\end{itemize}

\section{Autoría oficial}
\VerbatimInput{../../../AUTORIA-ORCID-DOI-OFICIAL.md}

\section{Documentación funcional completa}
\subsection{README}
\VerbatimInput{../README.md}
\subsection{Manual de usuario}
\VerbatimInput{../MANUAL-USUARIO.md}
\subsection{Manual técnico}
\VerbatimInput{../MANUAL-TECNICO.md}
\subsection{Configuración de paquete}
\VerbatimInput{../pyproject.toml}
\subsection{Cita académica}
\VerbatimInput{../CITATION.cff}
\subsection{Licencia}
\VerbatimInput{../LICENSE}
\subsection{Configuración de pruebas}
\VerbatimInput{../pytest.ini}

\section{Código fuente completo}
\subsection{Inicialización de paquete}
\VerbatimInput{../src/padic_grn_viz/__init__.py}
\subsection{CLI}
\VerbatimInput{../src/padic_grn_viz/cli.py}
\subsection{Motor gráfico}
\VerbatimInput{../src/padic_grn_viz/engine.py}

\section{Pruebas completas}
\subsection{Prueba smoke}
\VerbatimInput{../tests/test_smoke.py}
\subsection{Pruebas de motor}
\VerbatimInput{../tests/test_engine.py}

\section{Evidencia textual de validación}
\subsection{Log de ejecución de pruebas}
\VerbatimInput{validacion_viz_pytest.txt}
\subsection{Entrada de referencia para visualización}
\VerbatimInput{../../../registro-dossier/06_evidencia_corridas/core_scores_p2.json}
\subsection{Nota}
Las salidas principales de este software son archivos PNG. Por requisito de entrega en texto,
este dossier incluye el codigo fuente completo y la entrada textual de referencia utilizada para generar las figuras.

\end{document}
